% !TEX program = xelatex
% ============================================================
%  مناقشة الأطروحة — Dissertation Defense Template
%  نموذج عرض تقديمي لمناقشة الأطروحة (Beamer) بدعم العربية (RTL)
%  Compile with: xelatex main.tex (run twice)
% ============================================================
%  Features:
%    - Beamer 16:9 aspect ratio
%    - Arabic (RTL) via polyglossia (loads bidi internally)
%    - Title slide: candidate, advisor, committee, university, date
%    - Outline, motivation, literature, methodology, results
%    - Discussion, future work, acknowledgments, thank-you slide
%    - English terms wrapped in \LR{...}
% ============================================================

\documentclass[aspectratio=169]{beamer}

% ---- Theme (keep simple to avoid bidi conflicts) ----
% TODO: يمكن تغيير السمة إلى Madrid أو أخرى حسب الرغبة.
\usetheme{default}
\usecolortheme{default}
\setbeamertemplate{navigation symbols}{}
\setbeamertemplate{footline}[frame number]

% ---- Core packages ----
\usepackage{graphicx}
\usepackage{array}
\usepackage{booktabs}
\usepackage{tabularx}
\usepackage{amsmath}

% ---- Fonts (before polyglossia) ----
\usepackage[no-math]{fontspec}
\usepackage[quiet]{polyglossia}

% ---- Fonts ----
\setmainfont[Scale=1]{Amiri-Regular.ttf}
\newfontfamily\arabicfont[Script=Arabic,Scale=0.95]{Amiri-Regular.ttf}
\newfontfamily\englishfont[Scale=0.95]{TeX Gyre Termes}
\setmonofont{Amiri-Regular.ttf}
\newfontfamily\arabicfontsf{Amiri-Regular.ttf}

% ---- Languages (polyglossia loads bidi internally) ----
\setdefaultlanguage[calendar=gregorian,hijricorrection=1]{arabic}
\setotherlanguage[variant=british]{english}

% ---- Hyperref (beamer loads hyperref; configure only) ----
\hypersetup{
  pdftitle={مناقشة الأطروحة},
  pdfauthor={اسم المرشح}
}

% ---- Captions in Arabic ----
\addto\captionsarabic{%
  \renewcommand{\contentsname}{الفهرس}%
  \renewcommand{\figurename}{الشكل}%
  \renewcommand{\tablename}{الجدول}%
}

% ---- Enumerate label fix (bidi + beamer recursion workaround) ----
% Note: beamer + bidi can recurse on \labelenumi; using itemize avoids it.

% ---- Table column types ----
\newcolumntype{L}[1]{>{\raggedright\arraybackslash}p{#1}}
\newcolumntype{C}[1]{>{\centering\arraybackslash}p{#1}}

\begin{document}

% ============================================================
% TITLE SLIDE
% ============================================================
\begin{frame}[plain]
\centering
\vspace{0.5em}
{\Large\bfseries اسم الجامعة}\\[0.3em]
{\normalsize القسم --- اسم القسم}\\[1.2em]
{\LARGE\bfseries عنوان الأطروحة (\LR{Thesis Title})}\\[1.5em]
\begin{tabular}{rl}
\textbf{المرشح (\LR{Candidate}):} & اسم المرشح \\[0.3em]
\textbf{المشرف (\LR{Advisor}):} & د. اسم المشرف \\[0.3em]
\textbf{لجنة المناقشة (\LR{Committee}):} & د. اسم العضو الأول \\[0.2em]
& د. اسم العضو الثاني \\[0.2em]
& د. اسم العضو الثالث \\
\end{tabular}\\[1.2em]
{\normalsize \today}
\end{frame}

% ============================================================
% OUTLINE
% ============================================================
\begin{frame}{المخطط (\LR{Outline})}
% TODO: عدّل بنود المخطط حسب عرضك.
\begin{itemize}
    \item الدافع والأهمية (\LR{Motivation})
    \item الأدبيات السابقة (\LR{Literature})
    \item المنهجية (\LR{Methodology})
    \item النتائج (\LR{Results})
    \item المناقشة (\LR{Discussion})
    \item العمل المستقبلي (\LR{Future Work})
    \item الشكر والتقدير (\LR{Acknowledgments})
\end{itemize}
\end{frame}

% ============================================================
% MOTIVATION
% ============================================================
\begin{frame}{الدافع والأهمية (\LR{Motivation})}
% TODO: اكتب الدافع والأهمية هنا.
\begin{itemize}
    \item المشكلة البحثية وأهميتها.
    \item الفجوة المعرفية التي تسعى الأطروحة لسدها.
    \item المساهمة المتوقعة للمجال.
\end{itemize}
\end{frame}

% ============================================================
% LITERATURE
% ============================================================
\begin{frame}{الأدبيات السابقة (\LR{Literature Review})}
% TODO: اكتب مراجعة الأدبيات هنا.
\begin{itemize}
    \item الدراسة الأولى (\LR{Author et al., 2023}): ملخص الإسهام.
    \item الدراسة الثانية (\LR{Author et al., 2022}): ملخص الإسهام.
    \item الدراسة الثالثة (\LR{Author et al., 2021}): ملخص الإسهام.
\end{itemize}
\end{frame}

% ============================================================
% METHODOLOGY
% ============================================================
\begin{frame}{المنهجية (\LR{Methodology})}
% TODO: اكتب منهجية البحث هنا.
\begin{itemize}
    \item تصميم الدراسة (\LR{Study Design}).
    \item جمع البيانات (\LR{Data Collection}).
    \item التحليل (\LR{Analysis}).
\end{itemize}
\end{frame}

% ============================================================
% RESULTS
% ============================================================
\begin{frame}{النتائج (1/3) (\LR{Results})}
% TODO: اكتب النتائج الأولى هنا.
\begin{itemize}
    \item النتيجة الرئيسية الأولى.
    \item النتيجة الرئيسية الثانية.
\end{itemize}
\end{frame}

\begin{frame}{النتائج (2/3) (\LR{Results})}
% TODO: اكتب النتائج الثانية هنا.

\begin{table}[H]
\centering
\small
\begin{tabularx}{0.85\textwidth}{C{2.5cm} C{2.5cm} C{2.5cm}}
\toprule
\textbf{العينة} & \textbf{القيمة} & \textbf{الوحدة} \\
\midrule
1 & \LR{10.5} & \LR{mg/L} \\
2 & \LR{12.3} & \LR{mg/L} \\
\bottomrule
\end{tabularx}
\end{table}
\end{frame}

\begin{frame}{النتائج (3/3) (\LR{Results})}
% TODO: اكتب النتائج الثالثة هنا.
% \begin{figure}
% \centering
% \includegraphics[width=0.7\textwidth]{figure.png}
% \caption{عنوان الشكل}
% \end{figure}
\centering
\fbox{\parbox[c][2.5cm][c]{0.6\textwidth}{\centering مكان الشكل (\LR{Figure placeholder})}}
\end{frame}

% ============================================================
% DISCUSSION
% ============================================================
\begin{frame}{المناقشة (\LR{Discussion})}
% TODO: اكتب مناقشة النتائج هنا.
\begin{itemize}
    \item تفسير النتائج الرئيسية.
    \item المقارنة مع الأدبيات السابقة.
    \item حدود الدراسة (\LR{Limitations}).
\end{itemize}
\end{frame}

% ============================================================
% FUTURE WORK
% ============================================================
\begin{frame}{العمل المستقبلي (\LR{Future Work})}
% TODO: اكتب اتجاهات العمل المستقبلي هنا.
\begin{itemize}
    \item توسيع نطاق الدراسة.
    \item تجريب فرضيات جديدة.
    \item تطبيقات عملية مقترحة.
\end{itemize}
\end{frame}

% ============================================================
% ACKNOWLEDGMENTS
% ============================================================
\begin{frame}{الشكر والتقدير (\LR{Acknowledgments})}
% TODO: اكتب نص الشكر والتقدير هنا.
\begin{itemize}
    \item شكر للمشرف (\LR{Advisor}).
    \item شكر للجنة المناقشة (\LR{Committee}).
    \item شكر للجهة الممولة (\LR{Funding}).
    \item شكر للعائلة والأصدقاء.
\end{itemize}
\end{frame}

% ============================================================
% THANK YOU
% ============================================================
\begin{frame}[plain]
\centering
\vspace{2em}
{\Huge\bfseries شكراً لكم}\\[1em]
{\Large على الإصغاء}\\[2em]
{\normalsize الأسئلة (\LR{Questions})}
\end{frame}

\end{document}
